\chapter{自由电子的量子态与相对论运动}
\label{chap:feqo-electron-state}

实验电子束既不是无限延展的动量本征态，也不只有一个确定的到达时刻。
量子力学允许我们叠加不同动量；相对论又规定每个动量分量以不同能量演化。
在计算电子同光场交换能量之前，必须先把这两件事放进同一个可归一化的状态。

\section{自由电子为什么会遇到量子光学问题}
\label{chap:feqo-phenomena}

一束没有经过调制的高能电子，在能谱仪上通常只留下以入射能量为中心的一条峰。让同一束电子掠过被激光照亮的纳米结构，谱线却会裂成一串等间隔峰，相邻两峰相差一个光子能量 \(\hbar\omega\)。如果只把过程理解成“一次吸收或一次发射”，最多只能预期中心峰两旁各多一条峰；实验中出现许多边带，说明电子在一次相干相互作用中同时积累了多条交换路径的振幅。

强度继续增大时，某些边带还会消失，随后再次出现。这不是某条跃迁通道被永久关闭，而是到达同一末态的振幅重新相消又相长。因而真正需要演化的是电子态的复振幅，而不只是若干互不相干的跃迁概率。

相互作用结束以后，电子能谱可以保持不变，能量分量之间的相位却仍在自由传播中继续积累。不同能量分量在某些传播距离重新同相，电子概率密度便在时间轴上收缩成尖峰。若光场本身也需要量子化，电子每交换一个量子都会改变光场；没有记录光场时，对光自由度求偏迹又会降低电子相干性。材料响应、探测器分辨率和数据反演随后再逐层改变实验真正看到的计数。

这本书因此追踪一条具体链条：给定入射电子以及经典光场、量子光场或材料环境，先求相互作用后的联合态，再传播到探测面，最后把理想概率变成探测器计数。每一层都保留自己的参数；初态能散、环境退相干、探测器展宽和反演正则化不会被混成同一个“模糊”。

\section{边带是能量相差 \texorpdfstring{$\hbar\omega$}{hbar omega} 的散射通道}

自由电子没有预先存在的等间隔能级。它的精确真空色散是
\begin{equation}
E(\vb{p})=\sqrt{m^2c^4+c^2\vb{p}^2},
\label{eq:feqo-relativistic-dispersion-preview}
\end{equation}
其中 \(m\) 是电子静质量，\(c\) 是光速，\(\vb{p}\) 是三维动量。后文所谓“第 \(\ell\) 条边带”，是相互作用挑出的动量通道，不是电子原来就有的束缚能级。

设入射电子能量为 \(E_0=E(p_0)\)，与角频率 \(\omega\) 的场
交换 \(\ell\) 个能量量子后，其末能量为
\(E_\ell=E_0+\ell\hbar\omega\)。若电子沿同一方向继续运动，
末动量并不是简单的 \(p_0+\ell\hbar\omega/v_0\)；
精确关系是
\[
p_\ell=\frac1c\sqrt{(E_0+\ell\hbar\omega)^2-m^2c^4}.
\]
这里选择沿原方向传播的正根，并要求
\(E_0+\ell\hbar\omega\ge mc^2\)；
达不到这个阈值的标号不对应自由正能电子的传播末态。
对 \(|\ell\hbar\omega|\ll E_0-mc^2\) 作 Taylor 展开，
才有 \(p_\ell=p_0+\ell\hbar\omega/v_0
+O((\ell\hbar\omega)^2)\)，其中
\(v_0=(\dd E/\dd p)_{p_0}\)。这个一次近似将成为后文
“低反冲”模型的入口，而精确式仍用来核查近似范围。

一次近似的误差可以直接从精确式读出，而不必用“电子很快”代替判断。
在固定传播方向的正动量支上，\(E'(p)=v(p)>0\)，
\(E''(p)=1/(\gamma^3m)\)，其中
\(\gamma=E/(mc^2)\)。对反函数 \(p(E)\) 求导，
\(p'(E)=1/v\) 且
\(p''(E)=-E''(p)/v^3\)。因此若
\(\delta E=\ell\hbar\omega\)，则
\begin{equation}
p(E_0+\delta E)=p_0+\frac{\delta E}{v_0}
-\frac{(\delta E)^2}{2\gamma_0^3mv_0^3}
+O((\delta E)^3).
\label{eq:feqo-sideband-momentum-recoil}
\end{equation}
二阶项相对一次项的大小是
\(|\delta E|/(2\gamma_0^3mv_0^2)\)。
它随 \(|\ell|\) 增长；所以单个光子造成的反冲小，
不表示最外侧边带仍可用同一个线性动量步长。

数量级可以把这些符号落到电子显微镜的能区。
取动能 \(K_0=\SI{200}{\kilo\electronvolt}\)、
\(mc^2\simeq\SI{511}{\kilo\electronvolt}\)，则
\(\gamma_0=1+K_0/(mc^2)\simeq1.391\)，
\(v_0=c\sqrt{1-\gamma_0^{-2}}\simeq0.695c\)。
若每次交换 \(\hbar\omega=\SI{2}{\electronvolt}\)，
一次交换的线性动量步长
\(\hbar\omega/v_0\simeq
\SI{1.54e-27}{\kilogram\metre\per\second}\)，
仅为入射动量的约 \(5.8\times10^{-6}\)。
这说明单步动量改变量小，却不保证多边带或长距离传播时
累积的二阶相位也小；后文会分别检查这两种小参数。

实验电子还具有有限能散。沿前进方向取归一化动量包络
\(a(p)\)，满足 \(\int_0^\infty |a(p)|^2\,\dd p=1\)。
由于 \(E(p)\) 在 \(p>0\) 严格递增，变量替换
\(\dd p=\dd E/v(E)\) 给出归一化能量包络
\begin{equation}
b(E)=\frac{a(p(E))}{\sqrt{v(E)}},\qquad
\int_{mc^2}^{\infty}|b(E)|^2\,\dd E=1.
\label{eq:feqo-energy-envelope-normalization}
\end{equation}
这里的 \(v^{-1/2}\) 是连续态换标号时不可省去的
雅可比因子。若相互作用产生能量平移
\(b_\ell(E)=b(E-\ell\hbar\omega)\)，
不同标号对应的两个波包一般仍可重叠：
\(\langle b_\ell,b_{\ell'}\rangle
=\int b_\ell(E)^*b_{\ell'}(E)\,\dd E\)
未必为零。只有初始能宽远小于 \(\hbar\omega\)，
且仪器也能分辨这段间隔，谱图上的峰才可直接当作
彼此分离的边带；否则要先把振幅相加，再取模平方。
在能量下限附近，平移后的定义域也要与
\(E\ge mc^2\) 相交，不能把形式上的负能量尾部当作传播态。

边带的强度不能只靠能量守恒求出，还取决于多条路径的相位。
为准确表达一种常见结果，先定义第一类整数阶 Bessel 函数为
周期函数的 Fourier 系数：
\[
J_\ell(z)=\frac1{2\pi}\int_0^{2\pi}
\ee^{\ii(z\sin\theta-\ell\theta)}\,\dd\theta,
\qquad \ell\in\mathbb Z.
\]
对固定实 \(z\)，被展开的周期函数光滑；
对系数积分分部任意多次可得
\(|J_\ell(z)|=O(|\ell|^{-N})\)（任意固定 \(N\)），
所以 Fourier 级数绝对收敛。其系数恰为上式，故恒等式是
\(\ee^{\ii z\sin\theta}=\sum_\ell J_\ell(z)
\ee^{\ii\ell\theta}\)。由周期指数函数的正交性，两边的
模方积分相等，立刻有 \(\sum_\ell|J_\ell(z)|^2=1\)
（此处 \(z\) 为实数）。将积分变量作 \(\theta\mapsto-\theta\)
又知 \(J_\ell(z)\) 为实数，故也可写成
\(\sum_\ell J_\ell(z)^2=1\)。在低反冲、单向传播且外场可视为给定
经典场时，后文将从演化算符推出
\begin{equation}
P_\ell=J_\ell^2(2|\beta|),\qquad \ell\in\mathbb Z,
\label{eq:feqo-pinem-preview}
\end{equation}
这里 \(\beta\) 是沿电子轨迹积累的无量纲复耦合。
一个 \(\beta\) 同时决定整条能量梳，所以不同边带不是彼此
独立的按钮。上面的级数恒等式说明该理想模型的总概率确实为一。
\(\beta\) 的相位在单次能谱中可能看不见，却会在第二次
相互作用或自由传播中重新变成可测人口。

当电子经过线性材料环境时，环境对任意电流源的响应可以写成推迟 Green 张量 \(\mathbf G\)。能量损失与辐射都来自同一个响应核，但实验投影的通道不同：电子能量损失谱统计电子少了多少能量，阴极发光只统计进入收集光学系统的辐射通道。后文会从 Maxwell 方程推出 \(\mathbf G\)，而不是把它当作一个现成的抽象符号。

\section{从制备到计数的五个不同问题}

制备层给出电子波包、光场态和材料初态；相互作用层给出 哈密顿量 或经典响应；传播层积累自由相位并容纳噪声；测量层把末态变成有限分辨率的概率；推断层再从有限计数估计未知参数。理想理论若直接与实验柱状图比较，常见错误正是漏掉了其中一层。

下一节先处理最靠前的问题：在 \(\SI{10}{\kilo\electronvolt}\) 到 \(\SI{1}{\mega\electronvolt}\) 的能区，怎样描述一个既相对论又具有有限宽度的电子波包，以及它怎样与电磁势耦合。

\section{练习}

\begin{exercise}
用\cref{eq:feqo-relativistic-dispersion-preview} 求群速度 \(\nabla_{\vb{p}}E\)，并检查低速极限是否回到 \(\vb{p}/m\)。
\end{exercise}

\begin{exercise}
说明“初态能散”和“探测器能量分辨率”分别属于五层模型中的哪一层，并举例说明二者为什么不能用同一个卷积核随意互换。
\end{exercise}


\begin{definition}[纯态、射线与归一化]
Hilbert 空间是内积诱导的距离下完备的向量空间。
电子纯态由其中的非零向量 \(|\psi\rangle\) 表示，但相差非零复数的
向量代表同一物理射线。计算概率时选取归一化代表，
\[
\langle\psi|\psi\rangle=1.
\]
正交投影是满足 \(\Pi^2=\Pi=\Pi^\dagger\) 的算符。
对这样的 \(\Pi\)，Born 规则给事件概率
\[
P(\Pi)=\langle\psi|\Pi|\psi\rangle.
\]
\end{definition}

这个概率的范围不是额外假设。由 \(\Pi^2=\Pi=\Pi^\dagger\)
有 \(P(\Pi)=\langle\Pi\psi|\Pi\psi\rangle=\|\Pi\psi\|^2\ge0\)。
而 \(\Pi\psi\) 与 \((I-\Pi)\psi\) 正交，故
\(1=\|\psi\|^2=\|\Pi\psi\|^2+\|(I-\Pi)\psi\|^2\)，
从而 \(P(\Pi)\le1\)。测量“落在某些边带”时，
\(\Pi\) 就是这些正交边带态的投影之和。

整体相位不可观测，不等于相对相位不可观测。边带态
\(\sum_\ell c_\ell|\ell\rangle\) 的共同相位可以删去，\(c_\ell c_{\ell'}^*\)
中的相对相位却会在传播和第二次相互作用中变成人口差。

\begin{definition}[连续动量基与波包]
三维 Dirac delta 是由
\(\int\delta^{(3)}(\vb{p}-\vb{p}_0)f(\vb{p})\,\dd^3p
=f(\vb{p}_0)\) 对光滑测试函数 \(f\) 定义的分布，
不是普通的平方可积函数。采用
\[
\langle\vb{p},s|\vb{p}',s'\rangle
=\delta_{ss'}\delta^{(3)}(\vb{p}-\vb{p}')
\]
的归一化。可归一化单电子波包写成
\[
|\Psi\rangle=\sum_s\int\dd^3p\,
\psi_s(\vb{p})|\vb{p},s\rangle,
\qquad
\sum_s\int\dd^3p\,|\psi_s(\vb{p})|^2=1.
\]
\end{definition}

平面波 \(|\vb{p},s\rangle\) 本身不属于 Hilbert 空间，而是广义本征态；真正的
电子束由平方可积包络 \(\psi_s\) 描述。将上式代入内积并先对
\(\vb{p}'\) 使用 delta 分布，便得到
\(\langle\Psi|\Psi\rangle=\sum_s\int\dd^3p\,|\psi_s|^2=1\)。
因此积分测度中没有额外的 \((2\pi\hbar)^{-3}\)；
它已经包含在位置表象的 \(\langle\vb{r}|\vb{p}\rangle\) 中。
能散、束腰和脉冲长度都属于制备层参数。

\section{从相对论色散到 Dirac 方程}
\label{sec:feqo-dirac-from-dispersion}

量子力学的时间演化对时间是一阶的，因而初始波函数决定以后各时刻的状态。
相对论能量关系却是 \(E^2=m^2c^4+c^2|\mathbf p|^2\)。
直接把正平方根 \(E(\mathbf p)\) 变成微分算子，会得到空间中的非局域算子。
Dirac 的做法是扩大波函数的分量数，用矩阵把这个二次关系分解成
一阶的局域方程。

先写 \(H_0=c\alpha_i p_i+\beta mc^2\)，要求
\(H_0^2=(m^2c^4+c^2|\mathbf p|^2)I\) 对每个 \(\mathbf p\)
成立。分别比较 \(p_i^2\)、\(p_ip_j\) 和 \(p_i\) 的系数，必须有
\begin{equation}
\alpha_i\alpha_j+\alpha_j\alpha_i=2\delta_{ij}I,\qquad
\alpha_i\beta+\beta\alpha_i=0,\qquad \beta^2=I.
\label{eq:feqo-dirac-anticommutators}
\end{equation}
这里的反对易子定义为 \(\{A,B\}=AB+BA\)。
取量子力学中自旋 \(1/2\) 的三个 Pauli 矩阵
\[
\sigma_1=\begin{pmatrix}0&1\\1&0\end{pmatrix},\quad
\sigma_2=\begin{pmatrix}0&-\ii\\\ii&0\end{pmatrix},\quad
\sigma_3=\begin{pmatrix}1&0\\0&-1\end{pmatrix},
\]
它们满足 \(\sigma_i\sigma_j+\sigma_j\sigma_i=2\delta_{ij}I_2\)。
令
\[
\alpha_i=\begin{pmatrix}0&\sigma_i\\\sigma_i&0\end{pmatrix},
\qquad
\beta=\begin{pmatrix}I_2&0\\0&-I_2\end{pmatrix},
\]
逐块相乘便验证式~\eqref{eq:feqo-dirac-anticommutators}。
电子波函数因此是四分量复列向量，满足
\begin{equation}
\ii\hbar\partial_t\Psi
=H_0\Psi,\qquad
H_0=c\boldsymbol\alpha\cdot(-\ii\hbar\nabla)+\beta mc^2.
\label{eq:feqo-dirac-free}
\end{equation}
这并没有把电子的四个分量误当成四个独立粒子；
矩阵关系保证每个平面波仍满足原相对论色散。

\section{正能旋量从本征方程算出}
\label{sec:feqo-positive-spinor}

代入 \(\Psi=u(\mathbf p)
\ee^{\ii(\mathbf p\cdot\mathbf r-Et)/\hbar}\)。
写 \(u=(\phi,\chi)^\mathsf T\)，其中上下两块各有两个复分量，
则 \(H_0u=Eu\) 等价于
\[
(E-mc^2)\phi=c\boldsymbol\sigma\cdot\mathbf p\,\chi,\qquad
(E+mc^2)\chi=c\boldsymbol\sigma\cdot\mathbf p\,\phi.
\]
正能根为 \(E_{\mathbf p}=\sqrt{m^2c^4+c^2|\mathbf p|^2}>0\)。
第二式给 \(\chi=c(\boldsymbol\sigma\cdot\mathbf p)\phi/
(E_{\mathbf p}+mc^2)\)；代入第一式，利用
\((\boldsymbol\sigma\cdot\mathbf p)^2=|\mathbf p|^2I_2\)，
恰还原 \(E_{\mathbf p}^2-m^2c^4=c^2|\mathbf p|^2\)。

任取二分量单位正交基 \(\xi_s\)（\(s=1,2\)），令
\(\phi=N_{\mathbf p}\xi_s\)。归一化条件是
\[
1=u_s^\dagger u_s
=N_{\mathbf p}^2
\left(1+\frac{c^2|\mathbf p|^2}{(E_{\mathbf p}+mc^2)^2}\right)
=N_{\mathbf p}^2\frac{2E_{\mathbf p}}{E_{\mathbf p}+mc^2}.
\]
故
\begin{equation}
u_s(\mathbf p)=
\sqrt{\frac{E_{\mathbf p}+mc^2}{2E_{\mathbf p}}}
\begin{pmatrix}
\xi_s\\[1mm]
\dfrac{c\boldsymbol\sigma\cdot\mathbf p}{E_{\mathbf p}+mc^2}\xi_s
\end{pmatrix},
\qquad
u_s^\dagger u_{s'}=\delta_{ss'}.
\label{eq:feqo-positive-spinor}
\end{equation}
对负能根 \(-E_{\mathbf p}\)，先取下分量
\(\chi=N_{\mathbf p}\eta_s\)，第一块方程给
\(\phi=-c(\boldsymbol\sigma\cdot\mathbf p)\chi/(E_{\mathbf p}+mc^2)\)。
因此另一组归一化本征旋量为
\[
v_s(\mathbf p)=N_{\mathbf p}
\begin{pmatrix}
-\dfrac{c\boldsymbol\sigma\cdot\mathbf p}
{E_{\mathbf p}+mc^2}\eta_s\\[1mm]\eta_s
\end{pmatrix},
\qquad H_0v_s=-E_{\mathbf p}v_s .
\]
其中 \(\eta_s\) 为二分量正交基。
按上下块相乘可查 \(u_s^\dagger v_{s'}=0\)；
正负两组各有两个态，共同张成四维分量空间。
本书多数实验只制备正能电子，但负能子空间在判断投影近似是否有效时不能
从代数上删去。

\begin{proposition}[自由 哈密顿量 的正负能投影]
固定动量 \(\mathbf p\)，定义
\[
\Lambda_\pm(\mathbf p)
=\frac12\left(I_4\pm\frac{H_0(\mathbf p)}{E_{\mathbf p}}\right).
\]
则 \(\Lambda_\pm^2=\Lambda_\pm\)、
\(\Lambda_+\Lambda_-=0\)、\(\Lambda_++\Lambda_-=I_4\)，且
\(H_0\Lambda_\pm=\pm E_{\mathbf p}\Lambda_\pm\)。
\end{proposition}
\begin{proof}
先用式~\eqref{eq:feqo-dirac-anticommutators} 得
\(H_0(\mathbf p)^2=E_{\mathbf p}^2I_4\)。
于是
\((I\pm H_0/E_{\mathbf p})^2
=2(I\pm H_0/E_{\mathbf p})\)，
正是幂等性；两个不同号的乘积为
\(I-H_0^2/E_{\mathbf p}^2=0\)。
最后将 \(H_0\) 乘进定义并再次用平方恒等式。
\end{proof}

投影算符能立即做一件事：对任意四分量动量振幅
\(b(\mathbf p)\)，\(\Lambda_+b\) 提取正能部分，
\(\Lambda_-b\) 提取负能部分。若外场随时间变化，两个子空间一般会混合；
只有混合幅度相对能隙足够小，才能把后者作为高阶修正。

\section{概率守恒与波包归一化}
\label{sec:feqo-current-and-packet}

对式~\eqref{eq:feqo-dirac-free} 左乘 \(\Psi^\dagger\)，
再对其 厄米 共轭右乘 \(\Psi\)，质量项因 \(\beta^\dagger=\beta\)
相消；\(\alpha_i^\dagger=\alpha_i\) 使空间导数合成散度：
\[
\partial_t(\Psi^\dagger\Psi)
=-c\bigl[(\nabla\Psi^\dagger)\cdot\boldsymbol\alpha\Psi
+\Psi^\dagger\boldsymbol\alpha\cdot\nabla\Psi\bigr]
=-\nabla\cdot(c\Psi^\dagger\boldsymbol\alpha\Psi).
\]
因此定义
\begin{equation}
\rho=\Psi^\dagger\Psi,\qquad
\mathbf j=c\Psi^\dagger\boldsymbol\alpha\Psi,\qquad
\partial_t\rho+\nabla\cdot\mathbf j=0.
\label{eq:feqo-dirac-current}
\end{equation}
\(\rho\) 的单位是 \(\si{\per\cubic\metre}\)，
\(\mathbf j\) 是概率流而非电荷流；电子电荷流为 \(q\mathbf j\)。
若波函数在无穷远衰减，使边界通量为零，对全空间积分便得到
\(\dd\int\rho\,\dd^3r/\dd t=0\)。

对正能平面波，还能核对流速。把式~\eqref{eq:feqo-positive-spinor}
代入 \(u_s^\dagger\alpha_i u_s\)，上下块相乘后利用
\(\sigma_i(\boldsymbol\sigma\cdot\mathbf p)
+(\boldsymbol\sigma\cdot\mathbf p)\sigma_i=2p_iI_2\)，得到
\[
c\,u_s^\dagger\alpha_i u_s
=\frac{c^2p_i}{E_{\mathbf p}}
=\frac{\partial E_{\mathbf p}}{\partial p_i}.
\]
概率流除以概率密度，正是色散的群速度。

平面波不能在全空间归一化。全书取
\[
\langle\mathbf r|\mathbf p,s\rangle
=(2\pi\hbar)^{-3/2}u_s(\mathbf p)
\ee^{\ii\mathbf p\cdot\mathbf r/\hbar}.
\]
用 Fourier 恒等式
\(\int\dd^3r\,\ee^{\ii(\mathbf p'-\mathbf p)\cdot\mathbf r/\hbar}
=(2\pi\hbar)^3\delta^{(3)}(\mathbf p'-\mathbf p)\)，
并在 delta 支持处使用 \(u_s^\dagger u_{s'}=\delta_{ss'}\)，得
\(\langle\mathbf p,s|\mathbf p',s'\rangle
=\delta_{ss'}\delta^{(3)}(\mathbf p-\mathbf p')\)。
因此归一化的正能波包为
\begin{equation}
\Psi(\mathbf r,t)=
\sum_s\int\frac{\dd^3p}{(2\pi\hbar)^{3/2}}
a_s(\mathbf p)u_s(\mathbf p)
\ee^{\ii(\mathbf p\cdot\mathbf r-E_{\mathbf p}t)/\hbar},
\qquad
\sum_s\int\dd^3p\,|a_s(\mathbf p)|^2=1.
\label{eq:feqo-positive-wavepacket}
\end{equation}
初始动量宽度、能量宽度和脉冲长度由 \(a_s\) 决定，
不能在传播或探测阶段重新当作自由参数调节。

\section{窄带传播究竟舍弃了什么}
\label{sec:feqo-dispersion-expansion}

取中心动量 \(\mathbf p_0=p_0\hat{\mathbf z}\)，写
\(\delta\mathbf p=\mathbf p-\mathbf p_0\)、
\(E_0=E_{\mathbf p_0}\)、
\(\gamma_0=E_0/(mc^2)\)。
对精确色散求导，
\[
\partial_iE=\frac{c^2p_i}{E},\qquad
\partial_i\partial_jE
=\frac{c^2}{E}\delta_{ij}
-\frac{c^4p_ip_j}{E^3}.
\]
在 \(\mathbf p_0\) 处，纵向二阶导数为
\(1/(\gamma_0^3m)\)，两个横向二阶导数为
\(1/(\gamma_0m)\)，混合导数为零。因此
\begin{equation}
E(\mathbf p)
=E_0+v_0\delta p_z
+\frac{\delta p_\perp^2}{2\gamma_0m}
+\frac{\delta p_z^2}{2\gamma_0^3m}
+R_3(\delta\mathbf p).
\label{eq:feqo-dispersion-hessian}
\end{equation}
若波包支集位于中心周围半径 \(\Delta p\) 的球内，
且该球中三阶导数的算子范数不超过 \(M_3\)，Taylor 定理给
\(|R_3|\le M_3(\Delta p)^3/6\)。
传播时间为 \(t\) 时，由舍去余项造成的相位误差至多
\(tM_3(\Delta p)^3/(6\hbar)\)。
这个无量纲数远小于 \(1\) 时二阶包络才可靠；
“高能电子”本身不是省略三阶项的理由。

二阶近似能保留多少传播信息？用一维高斯电子包从头算一次。
固定一个正能自旋态并只看纵向。以 \(p_0\) 为中心的归一化
动量包络取为
\[
a(p)=(2\pi\sigma_p^2)^{-1/4}
\exp\!\left[-\frac{(p-p_0)^2}{4\sigma_p^2}\right],
\qquad \int_{-\infty}^{\infty}|a(p)|^2\,\dd p=1.
\]
假设 \(\sigma_p\ll p_0\)，负动量尾部才可忽略，
这使单向传播的读法与积分范围相容。
设 \(m_\parallel=\gamma_0^3m\)，
并在式~\eqref{eq:feqo-dispersion-hessian} 中只保留
\(E(p)=E_0+v_0\delta p+\delta p^2/(2m_\parallel)\)。
同时把随动量缓慢变化的旋量 \(u_s(p)\) 取为 \(u_s(p_0)\)；
在有效动量窗口中，这一步的幅度误差受
\(\sigma_p\sup\|\partial_pu_s\|\) 控制。
把此式与 \(a(p)\) 代入一维 Fourier 积分，令
\(\xi=z-v_0t\) 与 \(\sigma_{z0}=\hbar/(2\sigma_p)\)，
待积分的指数就是
\[
-\frac{\delta p^2}{4\sigma_p^2}
\left(1+\ii\frac{2\sigma_p^2t}{m_\parallel\hbar}\right)
+\frac{\ii\delta p\xi}{\hbar}.
\]
对 \(\delta p\) 配方并使用高斯积分，得到略去整体载波
\(\ee^{\ii(p_0z-E_0t)/\hbar}\) 后的包络
\[
\psi_{\rm env}(\xi,t)
=\frac{1}{(2\pi\sigma_{z0}^2)^{1/4}}
\frac{1}{\sqrt{1+\ii\tau}}
\exp\!\left[-\frac{\xi^2}{4\sigma_{z0}^2(1+\ii\tau)}\right],
\qquad
\tau=\frac{\hbar t}{2m_\parallel\sigma_{z0}^2}.
\]
取模平方时用
\(\operatorname{Re}(1+\ii\tau)^{-1}=(1+\tau^2)^{-1}\)，
得到可直接积分检验归一化的密度
\[
|\psi_{\rm env}(\xi,t)|^2
=\frac{\exp[-\xi^2/(2\sigma_{z0}^2(1+\tau^2))]}
{\sqrt{2\pi}\,\sigma_{z0}\sqrt{1+\tau^2}}.
\]
所以位置方差为
\[
\sigma_z^2(t)=\sigma_{z0}^2(1+\tau^2)
=\sigma_{z0}^2+
\frac{\hbar^2t^2}{4m_\parallel^2\sigma_{z0}^2}.
\]
中心按群速度移动，宽度则由纵向有效质量 \(\gamma_0^3m\)
控制。这个结果依赖窄带二阶近似和固定自旋；
高斯分布虽没有严格有限的支集，但在
\(|\delta p|\le L\sigma_p\) 内可用前述 Taylor 余项界，
外部概率由高斯尾积分直接控制。
若 \(tM_3(L\sigma_p)^3/(6\hbar)\) 不小，
或选定 \(L\) 后的尾概率不可忽略，
就应回到精确色散的积分，不能把上式外推到任意远的传播距离。

\section{电磁势怎样进入 哈密顿量}
\label{sec:feqo-minimal-coupling}

电子电荷写作 \(q=-e\)，其中 \(e>0\)。
经典带电粒子的 拉格朗日量 是
\[
L=-mc^2\sqrt{1-v^2/c^2}
+q\mathbf v\cdot\mathbf A(\mathbf r,t)-q\phi(\mathbf r,t).
\]
定义 Lorentz 因子
\(\gamma=(1-v^2/c^2)^{-1/2}\)。
由 \(\mathbf p=\partial L/\partial\mathbf v\) 得
\(\mathbf p=\gamma m\mathbf v+q\mathbf A\)：
正则动量与机械动量 \(\gamma m\mathbf v\) 相差 \(q\mathbf A\)。
Legendre 变换给
\[
H_{\rm cl}
=\sqrt{m^2c^4+c^2|\mathbf p-q\mathbf A|^2}+q\phi.
\]
把同一个机械动量替换用于一阶 Dirac 哈密顿量，得到
\begin{equation}
H_D=c\boldsymbol\alpha\cdot
(-\ii\hbar\nabla-q\mathbf A)+\beta mc^2+q\phi .
\label{eq:feqo-minimal-coupling}
\end{equation}
这也固定了符号：电子 \(q=-e\)，因此不能把
\(\mathbf p-q\mathbf A\) 随手写成 \(\mathbf p-e\mathbf A\)。

规范变换
\(\mathbf A'=\mathbf A+\nabla\chi\)、
\(\phi'=\phi-\partial_t\chi\) 后，令
\(\Psi'=\ee^{\ii q\chi/\hbar}\Psi\)。
直接对指数求导便有
\[
(-\ii\hbar\nabla-q\mathbf A')\Psi'
=\ee^{\ii q\chi/\hbar}
(-\ii\hbar\nabla-q\mathbf A)\Psi,
\]
时间导数与 \(q\phi'\) 中的 \(\partial_t\chi\) 同样相消。
所以同一概率流和实验预测不会随势的规范表示改变。

上面的连续性方程只写了同一波函数的概率守恒。
以后求跃迁振幅时出现的是两个不同电子态之间的电流，
因而还需要它的双线性版本。设 \(\Psi_i,\Psi_f\)
都是同一实电磁势下 Dirac 方程的解，定义
\[
\rho_{fi}^{\rm ch}=q\Psi_f^\dagger\Psi_i,\qquad
\mathbf j_{fi}^{\rm ch}=qc\Psi_f^\dagger
\boldsymbol\alpha\Psi_i .
\]
这两个量一般是复数，并不是一次实验中的实电荷密度或实电流；
它们是跃迁矩阵元。把 \(\Psi_i\) 的方程左乘
\(\Psi_f^\dagger\)，再从 \(\Psi_f\) 方程的共轭式减去，
同一标势与矢势项逐项抵消，得到
\begin{equation}
\partial_t\rho_{fi}^{\rm ch}
+\nabla\cdot\mathbf j_{fi}^{\rm ch}=0 .
\label{eq:feqo-transition-current-continuity}
\end{equation}
若两态是自由定态，时间因子为
\(\exp[-\ii(E_i-E_f)t/\hbar]\)。
记跃迁频率 \(\omega_{if}=(E_i-E_f)/\hbar\)，
则式~\eqref{eq:feqo-transition-current-continuity}
在这一频率分量上读作
\(\nabla\cdot\mathbf j_{fi}^{\rm ch}
=\ii\omega_{if}\rho_{fi}^{\rm ch}\)。
所以纵向电流和电荷密度不能在跃迁计算中独立指定。

这条关系也给规范不变性一个直接的检验。
对一个守恒的跃迁源，势的线性耦合为
\(H_{fi}(t)=\int(\rho_{fi}^{\rm ch}\phi
-\mathbf j_{fi}^{\rm ch}\cdot\mathbf A)\,\dd^3r\)。
规范变换使它改变
\[
\delta H_{fi}
=-\int\!\bigl(\rho_{fi}^{\rm ch}\partial_t\chi
+\mathbf j_{fi}^{\rm ch}\cdot\nabla\chi\bigr)\,\dd^3r
=-\frac{\dd}{\dd t}
\int\rho_{fi}^{\rm ch}\chi\,\dd^3r ,
\]
最后一步用式~\eqref{eq:feqo-transition-current-continuity}
并假定空间边界项消失。总时间导数只改变端点态的相位，
不会改变完整过程的跃迁概率。如果只保留
\(-\mathbf j\cdot\mathbf A\) 却在有纵向场时扔掉
\(\rho\phi\)，这个检验会失败；后文使用纯辐射规范时
必须说明为何可以这样选取。

外场若含有足够强的频率或空间变化，就可能把
\(\Lambda_+\) 子空间耦合到 \(\Lambda_-\)；
把单电子正能动力学当作封闭模型需要比较
\(\|\Lambda_-H_I\Lambda_+\|\) 与最小正负能隙，
并检查外场频谱是否跨越该能隙。强场对产生区域必须使用量子场描述，
不能把式~\eqref{eq:feqo-minimal-coupling} 的单粒子解释无限延伸。

\begin{exercise}
由式~\eqref{eq:feqo-positive-spinor} 算
\(u_s^\dagger\boldsymbol\alpha u_s\)，核对
\(\mathbf j/\rho=c^2\mathbf p/E_{\mathbf p}\)。
\end{exercise}

\begin{exercise}
设波包主要满足 \(|\delta p_z|\le\Delta p\)。
用本节的余项界写出传播时间 \(t\) 后二阶近似相位误差小于
\(0.1\) 的充分条件，并说明缩短 \(t\) 与缩小 \(\Delta p\)
各会带来什么改变。
\end{exercise}


波包的时间演化由其动量振幅和色散关系共同决定；
电磁势怎样改变它，则已经由最小耦合式固定。
接下来还缺少另一半输入：给定电流怎样激发电磁场，
以及场的一个模式何时必须视为可交换的量子。
